\section{Obserwowalność}

\subsection{Macierz obserwowności}

Macierz obserwowności przyjmuje postać:

\begin{equation}
\mathcal{O} = \begin{bmatrix} C \\ CA \\ CA^2 \end{bmatrix} = \begin{bmatrix} 1 & 0 & 0 \\ 0 & 1 & 0 \\ 0 & 0 & 1{,}177 \end{bmatrix}
\end{equation}

\subsection{Rząd macierzy obserwowności}

Obliczając wyznacznik macierzy obserwowalności:

\begin{equation}
\det(\mathcal{O}) = 1 \cdot \begin{vmatrix} 1 & 0 \\ 0 & 1{,}177 \end{vmatrix} = 1 \times 1{,}177 = 1{,}177
\end{equation}

Ponieważ $\det(\mathcal{O}) = 1{,}177 \neq 0$, macierz obserwowności ma pełny rząd 3. Układ jest w pełni obserwowalny.


